\section{Introduction}\label{sec:introduction}

Mixed-integer formulations of nonlinear functions usually approximate
the function's graph, for example by piecewise-linear models or by
discretized squares and products. A formulation can use many interval
selectors, a few integer variables with large ranges, continuous lifting
variables, or long coefficients. An efficient encoding of one chosen
partition gives an upper bound, but it does not show what is possible
with all convex mixed-integer formulations. We ask how many integer variables are
necessary for a prescribed accuracy, and we identify geometric and
algebraic quantities that govern this minimum.

For a continuous map $F:D\to\R^m$ and a prescribed error body $K$, an
admissible graph relaxation $R$ must satisfy
\[
 \{(x,F(x)):x\in D\}\subset R
 \subset\{(x,w):x\in D,\ w-F(x)\in K\}.
\]
Thus the relaxation retains every input-output pair and controls the
output error at the same input. We denote by $\pconv$ the least number
of general integer coordinates in a convex lift of any such $R$
(\cref{def:dimension}). Integer ranges, the finite continuous dimension
and real coefficients are unrestricted, and the lift need not be closed. The
comparison quantity $\pbin$ restricts the lift to linear constraints and
binary coordinates, so $\pconv\le\pbin$. When a polynomial-time
algorithm produces a formulation of polynomial rational encoding length,
we write $\pout$ for its binary count.

All results concern approximation of the entire graph on the stated
domain; restricting that domain by additional nonlinear equations can
change the lower bounds. Small integer dimension does not imply fast
optimization, an ideal linear relaxation or a short rational
description.

Which sets admit mixed-integer formulations at all is the subject of
representability theory. \citet{jeroslow1984} characterized the sets
representable by rational mixed-integer linear formulations, and
\citet{lubin2022} studied the mixed-integer convex analogue. On the
constructive side, logarithmic embeddings of polyhedral unions \citep{vielma2018}, shared
piecewise-linear formulations \citep{lyu2026}, minimum-breakpoint and
corridor approximation \citep{rebennack2015,codsi2025}, and compact
quadratic discretizations \citep{beach2022} give upper bounds. Our
accuracy-dependent question minimizes over all admissible relaxations $R$
and all convex lifts, including nonclosed ones. The results show when
such constructions use close to the fewest integer variables that the
geometry allows, even against general integers and arbitrary convex
constraints.

\subsection{Main contributions}

\paragraph{1. An exact precision law for quadratic systems.}
Let $F$ be a fixed real quadratic system on a full-dimensional bounded
box $B\subset\R^n$, and let $\mathcal H$ be the span of its Hessians.
With componentwise absolute error $\eps>0$, \cref{thm:ncrank} proves
\begin{equation}\label{eq:intro-ncrank}
 \pconv(\eps),\ \pbin(\eps)
 =\tfrac12\ncrank(\mathcal H)\log_2(1/\eps)+O_{F,B}(1)
 \qquad(\eps\downarrow0).
\end{equation}
Both minima have this expansion; they need not agree at a fixed
tolerance. Noncommutative rank is an established matrix-space invariant
\citep{ggow2020,iqs2018}. One characterization is
\[
 \ncrank(\mathcal H)
 =n-\max_{U\subseteq\R^n}
       \bigl(\dim U-\dim\mathcal H(U)\bigr),
 \qquad
 \mathcal H(U)=\operatorname{span}\{Hu:H\in\mathcal H,\ u\in U\},
\]
where the maximum is over linear subspaces. It records the largest loss
of dimension that all Hessians force on a common subspace.

The joint invariant is essential. For the cross product
$(u,v)\mapsto u\times v$ on six inputs, every scalar combination of the
three outputs has Hessian rank at most four, whereas the Hessian space
has noncommutative rank six. The coefficient in \eqref{eq:intro-ncrank}
is therefore three, which no single scalar combination detects.

For the lower bound, contacts with the exact graph whose integer
witnesses have the same parity have admissible midpoints; bounds on the
covariance and volume of these contacts limit how much of the domain one
parity class can cover. For the upper bound, a symmetric shrunk subspace
gives a coordinate decomposition with discretization depths proportional
to $\log_2(1/\eps)$, with coefficients in $\{0,1/2,1\}$ summing to
$\ncrank(\mathcal H)/2$; shared product relaxations then attain the
lower coefficient. Exact special cases are the square
(\cref{thm:square}), scalar quadratic rank (\cref{thm:scalar-rank}),
one-sided inertia for epigraphs and hypographs (\cref{thm:inertia}), and
fractional vertex covers for bilinear interaction graphs
(\cref{thm:graph-cover}). For rational data the coordinate change and
formulation are computable in polynomial time
(\cref{thm:rational-ncrank}). Local and global Hessian ranks, constant
scalar Hessian rank, and positive perspectives extend parts of the
analysis beyond quadratics
(\cref{thm:smooth-ranks,thm:constant-rank,prop:perspective}).

\paragraph{2. A finite-accuracy benchmark with a rational construction.}
An asymptotic coefficient does not describe accuracy thresholds caused
by small or highly unequal coefficients. For quadratic outputs with
Hessians $H_j$ on $[0,1]^n$ and positive component tolerances $\eps_j$,
define
\[
 D_* = \max\{\det P:0\preceq P\preceq I,
       \ \operatorname{tr}(H_jPH_jP)\le\eps_j^2\ \forall j\},
 \qquad \Phi=-\tfrac12\log_2 D_*.
\]
\Cref{thm:finite-covariance} places both minima within
$O(n\log(n+1))$ of $\Phi$, with constants independent of the
coefficients, tolerances and number of outputs; \cref{ex:covariance-gap}
shows that this benchmark can lose order $n\log n$. For rational data,
\cref{thm:rational-finite} constructs in polynomial time a rational MILP
of polynomial encoding length with $\pout\le\pconv+O(n\log(n+1))$. The
determinant constraints need not be Euclidean convex; the algorithm uses
their geodesic convexity, rational approximations and explicit
feasibility repairs. Further results give a posteriori certificates
(\cref{prop:covariance-certificate}), correlated and general output
error bodies (\cref{thm:grouped-covariance,thm:general-quadratic-body}),
a reduction to the common nonlinear input rank
(\cref{thm:input-quotient}), and sharper bounds for positive
semidefinite block, diagonal and integer-feature structure
(\cref{thm:block-psd,cor:diagonal-psd,thm:integer-features}). A
thin-domain example shows that a containing box can misrepresent a
correlated input domain (\cref{ex:thin-domain}). Finally,
\cref{thm:count-hardness} limits approximation of the minimum count:
for every fixed $\delta>0$, a polynomial-time algorithm that always
constructs a valid formulation within an additive $O(n^{1-\delta})$ of
the minimum, even for rational convex quadratic instances, would imply
$P=NP$.

\paragraph{3. Scalar convex functions with controlled descriptions.}
For every continuous convex scalar function on a compact interval and
every positive absolute tolerance, \cref{lem:scalar-chords} proves
\[
 \pconv\le\pbin\le\pconv+2.
\]
Thus general integer ranges and nonlinear convex constraints save at
most two integer coordinates in this setting. This bound uses real
coefficients and does not control continuous formulation size. For
nondecreasing curvature, an accuracy-dependent truncated curvature
measure $M_\eta$ determines both minima as $\log_2(1+M_\eta)$ up to
additive constants (\cref{thm:curvature-mass}). \Cref{thm:convex-hybrid}
constructs a rational MILP within eleven binaries of $\pconv$ for every
rational convex polynomial on $[0,1]$ in dense encoding and every
positive rational tolerance (seven binaries for monotone curvature,
\cref{thm:compiled-curvature}). Its running time and total encoding length are
polynomial in the input length, including the tolerance. This requires
compact access to the selected intervals even when their number is
exponential in the input length; certified curvature integration and
greedy segmentation provide it.
Related results cover separable sums (\cref{thm:separable-scalar}),
positive separable polynomial systems (nonnegative coefficients,
unconditional error bodies) with a double-logarithmic degree overhead,
including sparse input with exponents in binary
(\cref{thm:positive-loglog}), and, under allocation oracle assumptions,
powers with binary-encoded rational exponents
(\cref{thm:rational-powers}). For a power $x^q$ with $q>1$, pure
relative error near zero admits no lift with finitely many integer
variables, whereas on $[a,1]$ both minima have order $\log\log(1/a)$
(\cref{prop:relative-power}).

\paragraph{4. Coupled outputs.}
For several convex outputs of one input, shared breakpoints are already
available \citep{lyu2026}. With box tolerances, common partitions give
$\pout\le\pconv+11+\lceil\log_2m\rceil$ for $m$ dense convex polynomial
outputs (\cref{thm:convex-vector-overlay}), and overhead
$12+\lceil\log_2\sum_jD_j\rceil$ for arbitrary dense polynomial outputs,
summing the degrees $D_j\ge2$ of the nonaffine outputs
(\cref{thm:arbitrary-vector-overlay}). Let $r$ be the dimension of the
outputs' span modulo affine functions, called their
\emph{curvature rank}. For
$r\ge1$ and box error, \cref{thm:vector-rank} bounds the additional
binary count by $\lceil\log_2r\rceil$ plus an absolute constant, both
for real-coefficient lifts and for a polynomial rational construction
from dense polynomial data; affine outputs need no binaries. Selecting a
basis from the original convex outputs preserves the nonnegative chord
gaps needed for this bound. Positive facet scalarizations and positive
polar spanners extend the comparison to nonnegative facet bodies
(\cref{thm:vector-rank}) and to general unconditional error bodies, where
$\pbin\le\pconv+\lceil\log_2(8r^2-1)\rceil$
(\cref{thm:vector-oracle-rank}). Given a strong separation oracle for
the body and dense rational data, the compiled overhead is
$17+2\lceil\log_2r\rceil$, and the MILP uses an explicit rational
linear band instead of the oracle. For separable vector maps
with $n$ inputs, a common curvature basis gives overheads between
$n\lceil\log_2r\rceil+7n$ and $2n\lceil\log_2r\rceil+21n$, depending on
the error model and on whether a rational construction is required
(\cref{thm:separable-vector-rank}).

\paragraph{5. Separations and limits.}
Examples mark the limits of these comparisons. Fixed-degree convex
product graphs have exact counts $\pconv=n$ and
$\pbin=\lceil n\log_2 3\rceil$, so the binary/general-integer gap can
grow with input dimension (\cref{thm:exact-box-gap}). A nonconvex
scalar polynomial family has two general integers and growing binary
dimension (\cref{thm:nonconvex-gap}). A tilted error body transfers this
gap to componentwise convex outputs while its condition number grows
(\cref{prop:tilted-gap}). A family of convex powers shows that
scalarizations and uniform refinement do not determine the vector count
(\cref{prop:vector-power-obstruction}). For $x^{1/D}$ on $[0,1]$ at error
$1/16$, four binaries suffice, but every rational MILP has encoding
length $\Omega(D)$; for $D=2^B$ a rational mixed-integer second-order
cone formulation with four binaries has $O(B\log(B+2))$ bits
(\cref{thm:root-encoding,thm:root-conic-separation}). Integer count,
rational description length, input dimension and error geometry must
therefore be specified separately.

\subsection{Relationship to prior work}

\Cref{tab:result-map} summarizes the closest comparisons. Its last
column describes results proved here, not the novelty of all their
ingredients.
\begin{table}[htbp]
\centering\small
\begin{tabular}{@{}>{\raggedright\arraybackslash}p{.25\textwidth}>{\raggedright\arraybackslash}p{.32\textwidth}>{\raggedright\arraybackslash}p{.35\textwidth}@{}}
\toprule
Prior work & Established result & Development in this paper\\
\midrule
\citet{lubin2022} & Minimum MICP rank and midpoint lower bound & Accuracy-dependent contact-volume bounds for all intermediate graph relaxations\\
\citet{beach2022,ggow2020,iqs2018} & Compact quadratic approximations; noncommutative rank and capacity theory & Exact joint quadratic coefficient and uniform finite covariance comparison\\
\citet{rebennack2015,simchowitz2018,codsi2025} & Minimum-breakpoint models, chord estimates, and adaptive or greedy approximation & Two-bit scalar comparison and eleven-bit polynomial rational construction\\
\citet{lyu2026,awerbuch2004} & Common SOS2 encodings and barycentric spanners & Curvature-rank count comparisons with box, facet, and oracle error models\\
\bottomrule
\end{tabular}
\caption{Established methods and the complete graph-approximation comparisons developed here.}
\label{tab:result-map}
\end{table}

\paragraph{Representability and formulation size.}
\citet[Definition 4.3 and Lemma 4.1]{lubin2022} define minimum MICP rank
and show that $w$ pairwise midpoint-incompatible points require at least
$\lceil\log_2w\rceil$ integer variables. We use the same parity
mechanism and then control whole contact classes by covariance or chord
geometry. Taking closures of contacts, rather than of the original lift,
makes the volume arguments valid without measurable witness selections
or closed lifts. On the upper-bound side, \citet[Proposition 1]{vielma2018}
embeds finite unions of bounded rational polyhedra using distinct binary
codes; our elementary version, \cref{lem:disjunction}, also permits real
coefficients. This encodes the local graph bands once they are chosen;
it does not bound their number or the size of an implicit family.

Mixed-integer extension complexity also distinguishes integer count from
linear formulation size. \citet{cevallos2018} derive lower bounds on
integer count when the number of linear inequalities is restricted,
for formulations whose projected mixed-integer hull represents a
combinatorial polytope; \citet{schade2026} strengthen these bounds for
stable set and knapsack. Those size measures count inequalities rather
than rational encoding bits. Our model instead requires the projected set
itself to stay in the prescribed error band, and its minimum integer
count allows arbitrary convex lifts without a size restriction.

\paragraph{Quadratic discretization and matrix-space geometry.}
\citet{beach2022} give quadratic approximations whose binary, row and
continuous-variable counts grow logarithmically in reciprocal error.
Product envelopes \citep{mccormick1976}, multiparametric disaggregation
\citep{teles2013} and its refinements \citep{beach2024,beach2024b} are
related shared discretization tools. Logarithmic growth and binary
product linearization are therefore antecedents of our constructions;
our lower bounds show, for example, that the sawtooth square count is
optimal against arbitrary convex lifts (\cref{thm:square}).
\citet{ggow2020} develop operator-scaling and capacity theory, and
\citet{iqs2018} give constructive polynomial-time noncommutative-rank
certificates. We import these algebraic results; the symmetric
coordinate allocation and contact-volume arguments connect them to graph
accuracy.

Hessian-based anisotropic interpolation \citep{cao2007}, determinant
maximization under linear matrix inequalities \citep{vandenberghe1998},
geodesic optimization on positive definite matrices
\citep{sra2013,zhang2016} and convex-body separation and optimization
\citep{gls1981} are established tools. Our covariance constraints are
generally quadratic in $P$, and the comparison is with the minimum over
arbitrary convex lifts.

\paragraph{Adaptive approximation and coupled outputs.}
\citet{rebennack2015} compute minimum-breakpoint continuous
piecewise-linear approximations, including shifted under- and
overestimators; \citet[Lemma 2.1]{simchowitz2018} give the factor-two
chord/midpoint estimate; and \citet[Sections 3--4]{codsi2025} build
continuous corridor approximations from greedy maximal segments. These
works establish adaptive knots, chord
estimates and greedy segmentation. Our additional claims are the
comparison with arbitrary convex integer lifts and, for the polynomial
classes, control of the total rational encoding without enumerating all
segments. The compiler uses established continuous encodings of Boolean
circuits \citep{avis2019} and of products involving discrete functions
\citep{adams2012}.

\citet[Proposition 1]{lyu2026} share one SOS2 constraint among several
piecewise-linear functions of the same input. We build on this
possibility and compare the required common partition with $\pconv$
through curvature rank and error-body geometry. The determinant
selection is a barycentric-spanner method \citep{awerbuch2004};
approximate spanner construction \citep{plevrakis2020} and the
equivalence of separation and optimization \citep{gls1981} are also
inherited. For the separation examples, the cap-set estimate is due to
\citet{ellenberg2017}, and difference-of-convex decomposition is
classical \citep{hartman1959}.

\paragraph{Novelty and scope.}
To the best of our knowledge, the identification in
\eqref{eq:intro-ncrank} of half the noncommutative rank with the minimum
integer-dimension precision coefficient for quadratic graph approximation
is new. The algebraic invariant and its algorithms are not new; the
contribution is the lower bound for all admissible convex lifts together
with a matching binary linear construction. We do not claim optimality of every additive constant or a
characterization of all smooth and vector maps. In particular, whether
a universal additive binary/general-integer gap holds for one-input
convex vectors with growing output dimension and box error remains open
(\cref{subsec:one-input-open}).

\paragraph{Organization.}
\Cref{sec:foundations} defines the model and the contact arguments.
\Cref{sec:scalar-quadratic,sec:ncrank} prove the scalar and joint
quadratic precision laws, and \cref{sec:smooth} extends them to smooth
maps and perspectives. \Cref{sec:finite-quadratic,sec:rational-quadratic,sec:quadratic-bodies,sec:psd-structure}
treat finite accuracy, rational constructions, output error bodies and
positive structure; \cref{sec:quadratic-hardness} proves the hardness
result. \Cref{sec:scalar-geometry,sec:scalar-compilation} treat scalar
convex functions, \cref{sec:positive-allocation,sec:pure-powers} treat
positive separable polynomial systems and pure powers, and
\cref{sec:encoding-boundaries} treats relative error and the encoding
separation.
\Cref{sec:vector,sec:vector-oracle,sec:separable-vector} treat coupled
convex outputs; \cref{sec:vector-obstructions,sec:exact-separations}
give the separation examples.
Each section states its own hypotheses and defines the rank or
allocation benchmark it uses; these quantities are not interchangeable.

\paragraph{Formal verification scope.}
The accompanying Lean formalization, described in its README and
verification records, covers two groups of results. The first contains
the exact square count, the scalar quadratic rank and one-sided inertia
laws, their finite linear upper constructions, the exact convex one-sided
product count, and the continuous square-epigraph inequality lower bound.
The second contains the exact box counts of \cref{thm:exact-box-gap},
including closed lifts, the explicit rational formulation with $3n$
auxiliaries and $13n$ inequalities, and its monotone affine shear. Other
results, including the vector quadratic results, are not formally
verified.
